\section{Kontextfreie Sprachen}

\begin{frame}{Typ-2 Grammatiken: Normalformen}
	\begin{itemize}
		\item Vereinfachungen für kontextfreie Grammatiken $G$
		\begin{itemize}
			\item $\varepsilon$-Regeln eliminieren: $S \rightarrow aAb; A \rightarrow aS|\varepsilon$\\
			wird zu	$S \rightarrow aAb|ab; A \rightarrow aS$
			\item $G$ reduzieren: nutzlose Variablen eliminieren
			\item Kettenregeln $A \rightarrow B$ entfernen
		\end{itemize}
		\item Chomsky-Normalform
		\begin{itemize}
			\item zu jeder $\varepsilon$-freien, reduzierten kontextfreien Grammatik $G$ ohne Kettenregeln lässt sich eine äquivalente Grammatik $G'$ angeben, die nur Regeln der Art $A \rightarrow a$ bzw. $A \rightarrow BC$ enthält.
			\item Konstruktion: Ausgangspunkt $A\rightarrow X_1X_2 \ldots X_m, m\geq 2, X_i \in V \cup \Sigma$
			\begin{itemize}
				\item für $X_i=a\in \Sigma$ neue Variable $X_a$ mit $X_a \rightarrow a$ einführen
				\item verbleibende Regeln $A \rightarrow B_1B_2\ldots B_m, B_i \in V, m\geq 3$:\\
				ersetze $A \rightarrow B_1D_1, D_1\rightarrow B_2D_2, \ldots, D_{m-2}\rightarrow B_{m-1}B_m$ ($D_i$ neue Variablen)
			\end{itemize}
		\end{itemize}
		\item Greibach-Normalform: Regeln vom Typ $A \rightarrow aB_1B_2 \ldots B_m|a$
	\end{itemize}
\end{frame}

\begin{frame}{Beispiel: Erzeugung Chomsky-Normalform}
	$G=(\{S, A, B, C, D, E\}, \{a, b, c, d\}, P, S)$\\
	$P=\{$\\
	\qquad\qquad $S \rightarrow bA \mid aB \mid C \mid Cb,$\\
	\qquad\qquad $A \rightarrow bAA \mid aS \mid a \mid bE \mid bDb,$\\
	\qquad\qquad $B \rightarrow aBB \mid bS \mid b,$\\
	\qquad\qquad $C \rightarrow Cc \mid c,$\\
	\qquad\qquad $D \rightarrow d \mid \varepsilon $\\
	\qquad$\}$
\end{frame}

\begin{frame}{Worterkennung: Cocke-Younger-Kasami-Algorithmus (CYK)}
	polynomiell auch für nichtdeterministische Sprachen
	\begin{itemize}
		\item Ziel: Worterkennung und Konstruktion Ableitungsbaum bottom-up
		\item Ausgangspunkt: $G$ in Chomsky-Normalform, $w=a_1a_2 \ldots a_n \in \Sigma^*$
		\item definiere $V[i,j]:=\{A \in V \mid A \underset{G}{\Rightarrow}^* w_{ij} \}$ mit\\
		$w_{ij}=a_ia_{i+1}\ldots a_{i+j-1}, |w_{ij}|=j, 1\leq i \leq n, 1\leq j \leq n+1-i$
		\item es gilt: $w\in L$ gdw. $S \in V[1,n]$
		\item schrittweise Berechnung (Rückführung auf kürzere Teilworte)\\
		$j=1: V[i,1]=\{X \in V \mid (X\rightarrow a_i) \in P\}$\\
		$j>1: V[i,j]=\bigcup_{l=1}^{j-1}\{A\in V\mid \exists (A\rightarrow BC): B \in V[i,l] \land C \in V[i+l, j-l] \}$
		\item Wortproblem mit polynomiellem Zeitaufwand $O(n^3)$ lösbar
		\item Grammatik ist mehrdeutig, falls ein Eintrag $A \in V[i,j]$ auf mehr als eine Weise entsteht
	\end{itemize}
\end{frame}

\begin{frame}{Beispiel: CYK-Algorithmus}
	$$L=\{a^nb^n\mid n\geq 1\}$$
	Zugehörige Grammatik in Chomsky-Normalform:
	\begin{align*}
		G=(&\{A, B, C, S\}, \{a, b\}, P, S)\\
		P=\{&S \rightarrow AB,\\
		    &S \rightarrow AC,\\
		    &C\rightarrow SB,\\
		    &A\rightarrow a,\\
		    &B\rightarrow b\}
	\end{align*}
	
\end{frame}

\note{
	Betrachten des Wortes $w=aabb$\\
	\begin{columns}
		\column{.3\textwidth}
		\begin{itemize}
			\item[$(a)$] $V[1,1]=\{A\}$
			\item[$(a)$] $V[2,1]=\{A\}$
			\item[$(b)$] $V[3,1]=\{B\}$
			\item[$(b)$] $V[4,1]=\{B\}$
			\item[$(aa)$] $V[1,2]=\varnothing$
		\end{itemize}
		\column{.3\textwidth}
		\begin{itemize}
			\item[$(ab)$] $V[2,2]=\{S\}$
			\item[$(bb)$] $V[3,2]=\varnothing$
			\item[$(aab)$] $V[1,3]=\varnothing$
			\item[$(abb)$] $V[2,3]=\{C\} (l=2)$
			\item[$(aabb)$] $V[1,4]=\{S\} (l=1)$
		\end{itemize}
		\column{.4\textwidth}
		\centering
		\begin{tabular}{|c|c|c|c|}
			\hline
			$a$ & $a$ & $b$ & $b$ \\
			\hline
			\hline
			$\{A\}$ & $\{A\}$ & $\{B\}$ & $\{B\}$ \\
			\hline
			$\varnothing$ & $\{S\}$ & $\varnothing$ &  \\
			\hline
			$\varnothing$ & $\{C\}$ &  &  \\
			\hline
			$\{S\}$ &  &  &  \\
			\hline
		\end{tabular}
	\end{columns}
	Weiteres Beispiel: $S\rightarrow AB; A\rightarrow BA\mid a; B\rightarrow CC\mid b; C\rightarrow AB\mid c; w=caccab$
}

\subsection{Praktische Parser-Algorithmen}

\begin{frame}{Exkurs: LL- und LR-Parser}
	\begin{itemize}
		\item CYK: allgemein für jede kontextfreie Grammatik anwendbar, aber $O(n^3)$
		\item in der Praxis (Compilerbau): deterministische, lineare ($O(n)$) Parser für geeignete Teilklassen kontextfreier Grammatiken
		\item zwei Grundprinzipien
		\begin{itemize}
			\item \textbf{Top-Down (LL-Parser)}: Konstruktion des Ableitungsbaums von der Wurzel zu den Blättern, Linksableitung
			\item \textbf{Bottom-Up (LR-Parser)}: Konstruktion des Ableitungsbaums von den Blättern zur Wurzel, (umgekehrt rekonstruierte) Rechtsableitung
		\end{itemize}
		\item Namensgebung $LL(k)$/$LR(k)$: \underline{L}eft-to-right scan der Eingabe, \underline{L}inks- bzw. \underline{R}echtsableitung, $k$ Symbole Vorschau (Lookahead)
	\end{itemize}
\end{frame}

\begin{frame}{LL-Parser (Top-Down)}
	\begin{itemize}
		\item Prinzip: beginnend beim Startsymbol wird stets die am weitesten links stehende Variable ersetzt (Linksableitung); die passende Regel wird anhand der nächsten $k$ Eingabesymbole vorausschauend bestimmt
		\item Voraussetzung ($LL(k)$-Eigenschaft der Grammatik), u.a.
		\begin{itemize}
			\item keine Linksrekursion ($A \rightarrow A\alpha$)
			\item für Regel-Alternativen $A \rightarrow \alpha_1 \mid \alpha_2$: disjunkte $FIRST_k$-Mengen
		\end{itemize}
		\item Hilfsmittel: $FIRST$-Menge (mögliche Anfangssymbole einer Ableitung) und $FOLLOW$-Menge (mögliche Nachfolgesymbole einer Variable) $\rightarrow$ Parsertabelle $V \times \Sigma \rightarrow P$ (Variable $\times$ Vorschau $\rightarrow$ anzuwendende Regel)
		\item Realisierung: rekursiver Abstieg (\emph{recursive descent}, eine Prozedur je Variable) oder tabellengesteuert (Kellerautomat + Tabelle)
	\end{itemize}
\end{frame}

\begin{frame}{Beispiel: LL(1)-Parser}
	$G=(\{S,E,T,T',F\}, \{+,*,(,),id\}, P, S)$\\
	$P=\{S \rightarrow TE, \quad E \rightarrow {+}TE \mid \varepsilon, \quad T \rightarrow FT',$\\
	\qquad$T' \rightarrow {*}FT' \mid \varepsilon, \quad F \rightarrow (S) \mid id\}$\\[0.5em]
	Linksableitung von $w=id+id*id$ (stets am weitesten links stehende Variable ersetzt):\\
	\footnotesize
	$S \Rightarrow TE \Rightarrow FT'E \Rightarrow id\,T'E \Rightarrow id\,E \Rightarrow id{+}TE \Rightarrow id{+}FT'E$\\
	$\Rightarrow id{+}id\,T'E \Rightarrow id{+}id{*}FT'E \Rightarrow id{+}id{*}id\,T'E \Rightarrow id{+}id{*}id\,E \Rightarrow id{+}id{*}id$
\end{frame}

\begin{frame}{LR-Parser (Bottom-Up)}
	\begin{itemize}
		\item Prinzip: \emph{Shift-Reduce-Parsing} mit Stack
		\begin{itemize}
			\item \textbf{Shift}: nächstes Eingabesymbol auf den Stack legen
			\item \textbf{Reduce}: steht die rechte Seite einer Regel ("`Handle"') oben auf dem Stack, wird sie durch die linke Seite ersetzt
		\end{itemize}
		\item Ergebnis: (umgekehrt rekonstruierte) Rechtsableitung
		\item Varianten wachsender Mächtigkeit/Tabellengröße: $LR(0)$, $SLR(1)$, $LALR(1)$, $LR(1)$
		\item mächtiger als $LL(k)$: erkennt größere Klasse (deterministisch) kontextfreier Grammatiken, u.a. auch linksrekursive
		\item im Kern ein deterministischer Kellerautomat (DPDA), gesteuert über eine Zustands-/Aktionstabelle
	\end{itemize}
\end{frame}

\begin{frame}{Beispiel: LR-Parser (Shift-Reduce)}
	Grammatik: $S \rightarrow E, \quad E \rightarrow E+T \mid T, \quad T \rightarrow id$\\
	Eingabe $w=id+id$:\\[0.3em]
	\centering
	\begin{tabular}{|l|l|l|}
		\hline
		Stack & Rest-Eingabe & Aktion \\
		\hline
		\# & $id+id$ & shift \\
		\#$id$ & $+id$ & reduce $T\rightarrow id$ \\
		\#$T$ & $+id$ & reduce $E\rightarrow T$ \\
		\#$E$ & $+id$ & shift \\
		\#$E{+}$ & $id$ & shift \\
		\#$E{+}id$ &  & reduce $T\rightarrow id$ \\
		\#$E{+}T$ &  & reduce $E\rightarrow E{+}T$ \\
		\#$E$ &  & reduce $S\rightarrow E$ \\
		\#$S$ &  & akzeptiert \\
		\hline
	\end{tabular}
\end{frame}

\begin{frame}{LL- vs. LR-Parser: Vergleich}
	\begin{itemize}
		\item Ableitung: LL $\rightarrow$ Linksableitung (top-down), LR $\rightarrow$ Rechtsableitung, rückwärts rekonstruiert (bottom-up)
		\item Grammatikklassen: jede $LL(k)$-Grammatik ist auch $LR(k)$, umgekehrt gilt das nicht $\rightarrow$ LR ist strikt mächtiger
		\item Implementierung: LL oft handgeschrieben (rekursiver Abstieg), LR meist werkzeuggeneriert (Parsertabellen)
		\item Werkzeuge (Beispiele): ANTLR ($LL(*)$) vs. yacc/bison, CUP (LALR)
		\item gemeinsam: beide benötigen $O(n)$ Zeit (im Gegensatz zu CYK mit $O(n^3)$), dafür nur für Teilklassen kontextfreier Sprachen anwendbar
	\end{itemize}
\end{frame}

\begin{frame}{Pumping-Lemma für kontextfreie Sprachen}
	\begin{itemize}
		\item Pumping-Lemma: für jede kontextfreie Sprache $L$ gilt\\
		es existiert eine natürliche Zahl $n$ so, dass für alle $z \in L$ mit $|z|\geq n$ gilt:\\
		$z$ lässt sich zerlegen in $z=uvwxy$ mit
		\begin{itemize}
			\item $|vx|\geq 1$
			\item $|vwx| \leq n$
			\item $uv^iwx^iy \in L, i=0, 1, 2, \ldots$
		\end{itemize}
		\item Beweis:
		\begin{itemize}
			\item wähle $n=2^m$, $m$: Anzahl der Nichtterminalsymbole, wenn sich G in Chomsky-Normalform befindet
			\item mindestens eine Variable $A$ muss sich im Ableitungsbaum nach spätestens $m$ Schritten wiederholen;\\
			$A\underset{G}{\Rightarrow}^*\omega_1 A \omega_2$, wähle $v: \omega_1 \underset{G}{\Rightarrow}^* v, x: \omega_2 \underset{G}{\Rightarrow}^* x, w: A \underset{G}{\Rightarrow}^* w$;\\
			diese Ableitungen enthalten jeweils keine Wiederholungen von Variablen
		\end{itemize}
	\end{itemize}
\end{frame}

\note{
	Nachfolgendes Beispiel:\\
	$L=\left\{a^kb^kc^k\mid k \geq 1\right\}$
}

\begin{frame}{Beispiel: Pumping-Lemma für kontextfreie Sprachen}
	$L=\{a^kb^kc^k \mid k \geq 1 \}$
	\begin{itemize}
		\item $L$ ist kontext-sensitiv: Beweis Typ-1-Grammatik\\
		$G=(V, \Sigma, P, S), V=\{S, B, C\}$\\
		$P=\{S\rightarrow aSBC | aBC, CB \rightarrow BC,$\\
		\qquad $aB \rightarrow ab, bB \rightarrow bb, bC \rightarrow bc, cC \rightarrow cc \}$
		\item $L$ ist nicht kontextfrei: Beweis Pumping-Lemma\\
		wäre $L$ kontextfrei, sei $n$ das "`Pumping-$n$"'\\
		wähle $z=a^nb^nc^n \in L, |z|=3n \geq n$\\
		es ex. eine Zerlegung $z=uvwxy$ mit $|vx|\geq 1, |vwx| \leq n$\\
		$\rightarrow$ $vwx$ kann höchstens zwei verschiedene Buchstaben enthalten\\
		$\rightarrow$ in $uv^iwx^iy$ werden ein oder zwei Buchstaben "`aufgepumpt"', der dritte nicht\\
		$\rightarrow$ $uv^iwx^iy \notin L$ für $i \neq 1$ \Lightning 
	\end{itemize}
\end{frame}

\begin{frame}{Abschlusseigenschaften kontextfreier Sprachen}
	\begin{itemize}
		\item Kontextfreie Sprachen sind abgeschlossen unter
		\begin{itemize}
			\item Vereinigung: $L_1\cup L_2$\qquad $S \rightarrow S_1|S_2$
			\item Konkatenation: $L_1L_2$ \qquad $S\rightarrow S_1S_2$
			\item Stern-Produkt: $L^*$ \qquad $S_0\rightarrow S|\varepsilon, S\rightarrow S_1|S_1S$
			\item Spiegelung: ersetze in Chomsky-Normalform $A\rightarrow BC$ durch $A \rightarrow CB$
		\end{itemize}
		\item Kontextfreie Sprachen sind \underline{nicht} abgeschlossen unter
		\begin{itemize}
			\item [a)] Durchschnitt\\
			Gegenbeispiel: $L_1=\{a^kb^kc^l\mid k,l\geq 0\}, L_2=\{a^kb^lc^l\mid k,l\geq 0\}$ sind kontextfrei\\
			$L=L_1\cap L_2$ ist nicht kontextfrei
			\item [b)] Komplement\\
			$L_1 \cap L_2 = \overline{\overline{L_1} \cup \overline{L_2}}$: Abgeschlossenheit unter Vereinigung und Komplement würde zu Abgeschlossenheit unter Durchschnitt führen\\
			$\rightarrow$ Widerspruch zu a) \Lightning
			
		\end{itemize}
	\end{itemize}
\end{frame}

\begin{frame}{Abschlusseigenschaften deterministisch kontextfreier Sprachen}
	\begin{enumerate}[a)]
		\item Deterministisch kontextfreie Sprachen sind abgeschlossen unter Komplementbildung (beachte: Akzeptanz durch Endzustand!)
		\item Deterministisch kontextfreie Sprachen sind \underline{nicht} abgeschlossen unter
		\begin{enumerate}[i.]
			\item Durchschnitt $L_1 \cap L_2$
			\item Konkatenation $L_1L_2$
			\item Stern-Produkt $L^*$
			\item Vereinigung $L_1 \cup L_2$
		\end{enumerate}
		\item $L_1 \cap L_2$ ist allerdings det. kontextfrei (bzw. kontextfrei) falls $L_1$ regulär und $L_2$ det. kontextfrei (bzw. kontextfrei) ist
	\end{enumerate}
\end{frame}
